% TEMPORARY DRAFT. DO NOT \input THIS FILE INTO paper.tex YET.
% This fragment proposes replacement text for the Level-I conditional solver
% analysis.  It leaves the exact quotient in paper.tex unchanged.
% TODO(bibliography): add the final citation key for Furue and Kudo, PXL.
% The ell=2 arithmetic rows below were independently checked on 12 August 2026.

\section{A transverse Level-I slice}
\label{sec:pxl-level1-slice}

We now specialize the exact quotient to one Level-I parameter shape.  Its
parameters are $(v,o,\ell,r)=(48,16,2,2)$.  This is the only published shape
for which the solver estimate is close enough to the target that its constants
affect the headline conclusion.  We therefore state its exact reduction and
its conditional solver model separately.

Put $A=\F_{19^2}$ and identify $A^{64}$ with $\F_{19}^{128}$.  Write
$U=A^{64}=V_{\rm pub}\oplus O_{\rm pub}$, where
$V_{\rm pub}\cong A^{48}$ is the public-vinegar summand, and let
$\pi_v:U\to V_{\rm pub}$ be the projection.  Use the
zero-offset column relation from \Cref{eq:sq-l2-zero}.  After the public rank
test in \Cref{prop:sq-zero-split}, the verifier is equivalent to a target
consistency test followed by a quadratic residual map
\begin{equation}
 W_Rt=0,
 \qquad
 g(u)=z,
 \qquad
 g:A^{64}\longrightarrow\F_{19}^{48}.
 \label{eq:pxl-level1-residual}
\end{equation}
Here $z=G_Rt$ is uniform after conditioning on consistency.  The first test
passes with probability $19^{-16}$ for a uniform official target.  It is
performed before any nonlinear solve.

An injective linear map
\[
 T:\F_{19}^{h}\longrightarrow A^{64}
\]
defines the slice direction $L=\im T$.  Fix an $\F_{19}$-basis
$\ell_1,\ldots,\ell_h$ of $L$ and put
\begin{equation}
 L_A=\Span_A\{\pi_v(\ell_1),\ldots,\pi_v(\ell_h)\}
 \subseteq V_{\rm pub}.
 \label{eq:pxl-extension-span}
\end{equation}
The slice direction is \emph{$A$-independent} if
$\dim_A L_A=h$.  This condition is stronger than injectivity of
$\pi_v|_L$.  Its construction must also be independent of the fresh
coefficient region used by \Cref{lem:sq-full-domain-spectrum}.  An affine
coset $\mathcal A=a+L$ is \emph{$\ell=2$ transverse} if
\begin{equation}
 \pi_v(a)\notin L_A.
 \label{eq:pxl-l2-transverse}
\end{equation}
This condition is independent of the chosen representative of the coset.
The attack samples a uniform affine coset and solves all $48$ equations only
when \eqref{eq:pxl-l2-transverse} holds.

\begin{proposition}[Exact transverse affine restriction]
\label{prop:pxl-level1-transverse}
Assume that the supplied public key passes $\rank Q_R=48$ in
\Cref{prop:sq-zero-split}, and let $T$ be any injection of dimension
$h\le48$.  For every consistent target and every affine coset
$\mathcal A=a+\im T$, the complete substituted verifier is equivalent on
$\mathcal A$ to
\begin{equation}
 f_{a,z}(x):=g(a+Tx)-z=0,
 \qquad x\in\F_{19}^{h}.
 \label{eq:pxl-restricted-system}
\end{equation}
This is a system of $48$ affine quadratic equations in $h$ variables over
$\F_{19}$.  Every accepted root of \eqref{eq:pxl-restricted-system}
reconstructs a candidate that satisfies every equation in the original
verifier.  No residual equation is omitted.
\end{proposition}

\begin{proof}
The target split is an invertible public change of verifier outputs by
\Cref{prop:sq-zero-split}.  The map $x\mapsto a+Tx$ is an affine injection, so
substitution preserves the solution set inside $a+\im T$.  The reconstruction
map is the zero-offset instance of \eqref{eq:column-relation}.  The final
format check and the original verifier reject every invalid output.
\end{proof}

The PXL route uses $h=46=v-2$ and an $A$-independent direction.  It therefore
keeps every residual equation and produces two more equations than variables.
The expected root count of one uniform target and coset trial is $19^{-2}$.
This statement is an exact first-moment calculation.  It does not assert that
the restricted coefficients are uniform or that PXL terminates at the
predicted degree.

\subsection{Uniform target and coset sampling}

The root calculation uses a random affine coset, not a fixed affine
translate selected after inspecting the polynomial coefficients.  The next
lemma records the exact sampling rule.

\begin{lemma}[Uniform coset sampling]
\label{lem:pxl-uniform-coset}
Let $U$ be a finite-dimensional $\F_q$-space and let $L\subseteq U$ have
dimension $h$.  If $a$ is uniform in $U$, then $a+L$ is uniform in $U/L$.
Equivalently, choose any complement $U=L\oplus V$ and sample $a$ uniformly
in $V$.  These two procedures give the same distribution on affine cosets.
For every $x\in U$,
\begin{equation}
 \Pr[x\in a+L]=q^{h-\dim U}.
 \label{eq:pxl-point-in-coset}
\end{equation}
If the residual target $z$ is sampled independently and uniformly in
$\F_q^K$, then
\begin{equation}
 \Pr[x\in a+L\ \text{and}\ g(x)=z]
 =q^{h-\dim U-K}.
 \label{eq:pxl-point-root}
\end{equation}
Suppose in addition that $U=A^{v+o}$, that $L$ is $A$-independent in the
sense of \eqref{eq:pxl-extension-span}, and that $a$ is uniform in $U$.
Then
\begin{equation}
 \Pr[\pi_v(a)\notin L_A]=1-|A|^{h-v}.
 \label{eq:pxl-transverse-probability}
\end{equation}
\end{lemma}

\begin{proof}
Every coset has exactly $q^h$ representatives.  A uniform point of $U$
therefore induces the uniform distribution on $U/L$.  A fixed point belongs
to one of the $q^{\dim U-h}$ cosets, which proves
\eqref{eq:pxl-point-in-coset}.  Independence and uniformity of $z$ contribute
the factor $q^{-K}$ in \eqref{eq:pxl-point-root}.  The projection
$\pi_v(a)$ is uniform in $A^v$.  Since $L_A$ has $|A|^h$ elements,
\eqref{eq:pxl-transverse-probability} follows.
\end{proof}

The official target is made exactly uniform by
\eqref{eq:target-rejection-sampling}.  The output split is invertible, so the
consistency coordinates and the residual target are independent and uniform.
The attack first retains a consistent target and then samples the affine
coset independently.  Thus the experiment in \Cref{lem:pxl-uniform-coset}
is exact for the residual system.  The attack does not sample directly from
the transverse cosets.  It samples a uniform coset and discards a
nontransverse one.  This preserves the sampling experiment used by the root
theorem below.

\begin{lemma}[Mixed block support]
\label{lem:pxl-l2-mixed-block}
Let $C:A\times A\to\F_{19}$ be an $\F_{19}$-bilinear form and put
\[
 C^{\rm sym}(\xi,\zeta)=C(\xi,\zeta)+C(\zeta,\xi).
\]
The map $C\mapsto C^{\rm sym}$ onto the symmetric bilinear forms on the
two-dimensional $\F_{19}$-space $A$ is surjective.  Its kernel is the
one-dimensional space of alternating forms on $A$.
\end{lemma}

\begin{proof}
Since the characteristic is odd, a symmetric form $S$ has preimage $S/2$.
The kernel consists of the forms with
$C(\xi,\zeta)=-C(\zeta,\xi)$.  These are the alternating forms.  Their
space has dimension $\binom{2}{2}=1$.
\end{proof}

\begin{proposition}[$\ell=2$ affine coefficient support]
\label{prop:pxl-l2-affine-support}
Let $h\le v-1$, let $L$ be $A$-independent, and let $a+L$ be $\ell=2$
transverse.  For one fresh UOV base form, the restriction map from its fresh
native symmetric coefficients to the three descended $\ell=2$ channel
polynomials on $a+L$ is surjective onto
\begin{equation}
 \F_{19}[x_1,\ldots,x_h]_{\le2}^{\,3}.
 \label{eq:pxl-l2-affine-surjection}
\end{equation}
Across the $m$ independent base forms, the corresponding map is surjective
onto the coefficient space of $K=3m$ ordinary affine quadratic equations in
$h$ variables.  In particular, $h=v-2$ leaves transversality failure
probability $|A|^{-2}=19^{-4}$.
\end{proposition}

\begin{proof}
Put $b=\pi_v(a)$ and $e_i=\pi_v(\ell_i)$.  The two transversality conditions
give the direct sum
\[
 Ab\oplus Ae_1\oplus\cdots\oplus Ae_h\subseteq V_{\rm pub}.
\]
Let $B$ be one fresh native symmetric $\F_{19}$-bilinear form on this sum.
The diagonal block on $Ae_i$ is an arbitrary symmetric form on the
two-dimensional $\F_{19}$-space $A$.  It supplies the three channel
coefficients of $x_i^2$.  The block between $Ae_i$ and $Ae_j$ is an
arbitrary bilinear form on $A\times A$.  Its contribution to $x_ix_j$ is its
symmetrization.  By \Cref{lem:pxl-l2-mixed-block}, this contribution supplies
all three channel coefficients.  The only unused direction in the block is
its one-dimensional alternating kernel.

The block between $Ab$ and $Ae_i$ uses the same mixed block map and supplies
all three coefficients of $x_i$.  The diagonal block on $Ab$ supplies the
three constant coefficients.  The blocks are independent because the sum of
$A$-lines is direct.  Set every unused block to zero and extend $B$
arbitrarily to the rest of $V_{\rm pub}$.  This constructs a native preimage
for every triple of affine quadratic polynomials.

The block map is defined over $\F_{19}$.  One may verify surjectivity after
extending scalars to a splitting field.  If $C$ is the cokernel, scalar
extension gives $C\otimes_{\F_{19}}\Omega=0$.  Faithful flatness then gives
$C=0$.  Hence the split calculation descends to the official
$\F_{19}$ coefficient space.  Finally, the native coefficient blocks for
the $m$ base forms are independent parameters, so taking their direct sum
gives the stated map for $K=3m$ equations.  The incidence probability is
\eqref{eq:pxl-transverse-probability} with $h=v-2$.
\end{proof}

Surjectivity is a support statement.  It says that every ordinary affine MQ
coefficient tuple has a preimage.  If the fresh native symbols were uniform
over $\F_{19}$, the induced tuple would also be uniform because the map is
linear and surjective.  In the idealized XOF transcript used by the spectrum
theorem, the native symbols are independent byte reductions and are not
uniform over $\F_{19}$.  Full support therefore does not imply uniform output
coefficients, semi-regularity, or the PXL rank predictions.

\section{Accepted roots on a transverse slice}
\label{sec:pxl-root-probability}

This section keeps the collision argument independent of the solver.  It
applies to square and overdetermined slices and requires only $h\le K$ for
the full column rank conclusion.

\begin{theorem}[Root-collision identity]
\label{thm:collision-identity}
Let $U$ be an $\F_q$-space, let $L\subseteq U$ have dimension $h$, and let
$g:U\to\F_q^K$ have degree at most two.  Sample a uniform affine coset
$\mathcal A$ of $L$ in $U$ and an independent uniform target
$z\in\F_q^K$.  Put
\[
 Z=|\{x\in\mathcal A:g(x)=z\}|,
 \qquad \mu=q^{h-K}.
\]
For $0\ne y\in L$, define
\[
 T_y(x)=g(x+y)-g(x)-g(y)+g(0),
\]
and let $\kappa(y)$ be one if
$-(g(y)-g(0))\in\im T_y$ and zero otherwise.  Then
\begin{align}
 \mathbb EZ&=\mu,                                      \label{eq:pxl-moment-first}\\
 \mathbb E[Z(Z-1)]
 &=\mu\sum_{0\ne y\in L}\kappa(y)q^{-\rank T_y}
 \le\mu\sum_{0\ne y\in L}q^{-\rank T_y}.             \label{eq:pxl-moment-second}
\end{align}
The derivative $T_y$ has domain $U$ because the experiment averages over the
affine coset.
\end{theorem}

\begin{proof}
The proof is the counting argument already given in the current manuscript.
By \Cref{lem:pxl-uniform-coset}, every point of $U$ lies in the sampled coset
with probability $q^{h-\dim U}$ and matches the independent target with
probability $q^{-K}$.  This proves \eqref{eq:pxl-moment-first}.  For a pair of
distinct roots, write the second point as $x+y$ with $0\ne y\in L$.  The
collision condition is
\[
 T_y(x)=-(g(y)-g(0)).
\]
It has $\kappa(y)q^{\dim U-\rank T_y}$ solutions in $x$.  Dividing by the
number of target and coset pairs proves \eqref{eq:pxl-moment-second}.
\end{proof}

Let $X\subseteq U$ be the assignments accepted by the original signature
format rule and put $\alpha=|X|/|U|$.  Define the full-domain spectrum
\begin{equation}
 \bar\eta_L^U=\sum_{0\ne y\in L}q^{-\rank T_y}.
 \label{eq:pxl-eta}
\end{equation}
For a root $x$, let $D_Lg(x):L\to\F_q^K$ be the Jacobian restricted to the
slice directions.

\begin{theorem}[Accepted root with full column rank]
\label{thm:sq-accepted-root}
Assume $h\le K$, $\bar\eta_L^U\le\eta$, and
$\alpha>\eta/(q-1)$.  A uniform target and affine-coset trial contains an
accepted root $x$ satisfying
\[
 \rank D_Lg(x)=h
\]
with probability at least
\begin{equation}
 p_{\rm reg}(h,K,\alpha,\eta)
 =q^{h-K}\frac{a^2}{a+\eta},
 \qquad
 a=\alpha-\frac{\eta}{q-1}.
 \label{eq:pxl-preg}
\end{equation}
This bound does not assume independence between acceptance and Jacobian rank.
It also does not assume independence from the target or the quadratic map.
\end{theorem}

\begin{proof}
For each projective direction $[y]\in\mathbb P(L)(\F_q)$, consider the
condition
\[
 D_Lg(x)[y]=0.
\]
This is an affine system in $x$ with linear part $T_y$.  A union
bound over projective directions shows that the fraction of points where
$D_Lg(x)$ lacks full column rank is at most $\eta/(q-1)$.  If $Y$ counts the
accepted roots with full column rank in the sampled trial, then
\[
 \mathbb EY\ge q^{h-K}a.
\]
By \Cref{thm:collision-identity},
$\mathbb E[Y(Y-1)]\le q^{h-K}\eta$.  Applying the Cauchy Schwarz inequality
to $Y$ gives
\eqref{eq:pxl-preg}.
\end{proof}

For the Level-I zero-offset relation, a block is symmetric exactly when its
second free coordinate is zero.  The exact accepted density on $A^{64}$ is
therefore
\begin{equation}
 \alpha_{64}=19^{-64}\sum_{j=0}^{16}\binom{64}{j}18^{64-j}.
 \label{eq:pxl-alpha64}
\end{equation}
For a $46$-dimensional $A$-independent slice, set
\begin{equation}
 \eta_{\rm L1}=\eta_{128}(46,48)
 =(19^{46}-1)19^{-48}+2^{-128}.
 \label{eq:pxl-eta-level1}
\end{equation}
On the event $\bar\eta_L^U\le\eta_{\rm L1}$, the exact lower bound used in
the retry calculation is
\begin{equation}
 p_{\rm reg}^{\rm L1}
 =19^{-2}\frac{a_{\rm L1}^2}{a_{\rm L1}+\eta_{\rm L1}},
 \qquad
 a_{\rm L1}=\alpha_{64}-\frac{\eta_{\rm L1}}{18}.
 \label{eq:pxl-preg-level1}
\end{equation}
The attack also requires a transverse coset.  Since a uniform coset is
nontransverse with probability $19^{-4}$, a union bound gives the certified
lower bound
\begin{equation}
 p_{\rm dir}^{\rm L1}
 =\max\{p_{\rm reg}^{\rm L1}-19^{-4},0\}.
 \label{eq:pxl-pdir-level1}
\end{equation}
This subtraction does not assume independence between transversality and root
existence.  Thus $p_{\rm dir}^{\rm L1}$ is not claimed to equal the usable
trial probability.  It is a lower bound for the joint event consisting of
root existence, transversality, acceptance, and full column rank.  It makes no
claim about the behavior of PXL.

\section{Conditional PXL recovery}
\label{sec:pxl-conditional-recovery}

PXL first eliminates part of a Macaulay matrix over a polynomial ring in the
variables that will later be guessed.  It then evaluates those variables and
performs the remaining linear algebra over the base field.  Its published
cost analysis uses assumptions about the Hilbert series and the ranks of the
Macaulay matrices~\cite{FurueKudo2024PXL}.  We state those published model
assumptions first.  We then state a separate premise for applying the model
and the manuscript's root extraction step to each actual quotient instance.

\begin{definition}[Ordinary MQ regularity hypothesis $\mathsf H_{\rm PXL}$]
\label{def:pxl-regularity}
Let $K$ affine quadratic equations in $h$ variables be given.  Choose $k$
variables for the polynomial preprocessing step and let the remaining $h-k$
variables be the main variables.  For a degree bound $D$, the hypothesis
$\mathsf H_{\rm PXL}(K,h,k,D)$ consists of the following statements.
\begin{enumerate}[label=(\roman*)]
\item For every assignment to the $k$ guessed variables, the specialized
      system has at most one root over the algebraic closure, counted with
      multiplicity, and its highest-degree parts are semi-regular.
\item The homogeneous quadratic parts in the $h-k$ main variables are
      semi-regular, so their Hilbert series gives the stated residual matrix
      size.
\item The rows selected in each Macaulay degree have the predicted rank, and
      the matrix left after polynomial preprocessing has the predicted square
      size.
\item Degree $D$ suffices for the second linearization step to produce the
      elimination data predicted by the PXL analysis.
\end{enumerate}
These are ordinary MQ heuristics.  They are the affine semi-regularity,
projected semi-regularity, and Macaulay rank premises used in the published
PXL analysis.  They are not implied by coefficient support, and they are not
proved for the quotient systems in this paper.
\end{definition}

\begin{definition}[Instancewise route premise $\mathsf H_{\rm route}$]
\label{def:pxl-route}
Fix a public key, a retained target, and a transverse affine slice.  Let
$f_{a,z}$ be the resulting system in \eqref{eq:pxl-restricted-system}.  The
premise $\mathsf H_{\rm route}(f_{a,z};k,D)$ states the following facts about
this particular system.
\begin{enumerate}[label=(\roman*)]
\item Its PXL matrices have the degrees, dimensions, and ranks specified by
      $\mathsf H_{\rm PXL}(K,h,k,D)$.
\item If $f_{a,z}$ has an accepted base-field root with full column rank, the
      manuscript's specialization and root extraction procedure returns at
      least one such root from the PXL elimination data.
\end{enumerate}
The second item is an added solver completeness premise.  It is not part of
the published PXL assumptions.  The five PXL terms below account for the
published algebraic model.  The separate term $C_{\rm aux}$ retains the cost
of root extraction and every other operation not present in those five terms.
This premise is instancewise.  It makes no claim that the quotient systems
have the uniform MQ distribution, and it introduces no mixing probability.
\end{definition}

For fixed $(K,h,k,D)$, define the predicted residual matrix size by
\begin{equation}
 A=A(K,h,k,D)
 =\sum_{d=0}^{D}
   \max\left\{
     [t^d](1-t)^{K-h+k}(1+t)^K,
     0
   \right\}.
 \label{eq:pxl-A}
\end{equation}
The standard degree prediction chooses the least $D\ge2$ such that
\begin{equation}
 [t^D](1-t)^{K-h+k-1}(1+t)^K\le1.
 \label{eq:pxl-degree}
\end{equation}
Equation \eqref{eq:pxl-degree} is part of the conditional model.  It does not
prove a solving degree guarantee for the structured quotient system.

The PXL field-operation expression has five charged terms.  Put
$B=\binom{h-k+D}{D}$.  The two scalar preprocessing terms are
\begin{align}
 C_1
 &=B^\omega,                                           \label{eq:pxl-C1}\\
 C_2
 &=k^2(h-k)B^2.                                        \label{eq:pxl-C2}
\end{align}
The polynomial preprocessing term that clears the later Macaulay columns is
\begin{equation}
 C_3
 =k^2AB\binom{h+D}{D}.                                 \label{eq:pxl-C3}
\end{equation}
After preprocessing, evaluation and the second linearization cost
\begin{align}
 C_{\rm fix}
 &=19^kA^2\binom{k+D}{D},                             \label{eq:pxl-Cfix}\\
 C_{\rm lin}
 &=19^kA^\omega.                                      \label{eq:pxl-Clin}
\end{align}
The first three terms are paid once.  The last two include all $19^k$
assignments to the guessed variables.  Keeping only the rough dominant-term
formula would omit valid work from the published PXL model.

\begin{theorem}[Conditional Level-I PXL model]
\label{thm:pxl-level1-recovery}
Consider the Level-I shape $(48,16,2,2)$.  Assume the exact public rank check
of \Cref{prop:sq-zero-split}.  Assume the spectrum condition
$\bar\eta_L^U\le\eta_{\rm L1}$ for every invoked slice.  Finally, assume
that $\mathsf H_{\rm route}(f_{a,z};k,D)$ holds for every invoked restricted
system.  In this premise, the actual matrices satisfy the predictions of
$\mathsf H_{\rm PXL}(48,46,k,D)$ for the chosen parameters.  Charge every
$\F_{19}$ operation at $150$ Boolean gates.  This factor is a constructive
bound for one field operation.  It does not turn the PXL model into a bound
for the whole computation.
Put
\begin{equation}
 C_{\rm PXL}=C_1+C_2+C_3+C_{\rm fix}+C_{\rm lin},
 \qquad
 C_{\rm target}=2^{32}19^{16}.
 \label{eq:pxl-kernel-cost}
\end{equation}
The target term follows the manuscript convention of charging $2^{32}$ gates
for each retained uniform target and an expected $19^{16}$ targets for the
consistency test.  The target cost is paid anew in every root-theorem trial.
The unit-leading asymptotic estimate is
\begin{equation}
 W_{\rm est}^{\rm L1}
 =\frac{150C_{\rm PXL}+C_{\rm target}}
        {p_{\rm dir}^{\rm L1}}.
 \label{eq:pxl-estimated-gates}
\end{equation}
This formula does not charge setup, root extraction, control, storage traffic,
or final verification.  If $C_{\rm aux}$ accounts for these costs for one
trial, the corresponding augmented expression is
\begin{equation}
 W_{\rm aug}^{\rm L1}
 =\frac{150C_{\rm PXL}+C_{\rm target}+C_{\rm aux}}
        {p_{\rm dir}^{\rm L1}}.
 \label{eq:pxl-total-gates}
\end{equation}
These expressions account only for the named terms.  They leave the auxiliary
operations above unmeasured.

For the independently checked minimizing parameters, the Level-I model gives
the following two rows.  The parameter search is performed separately for
each value of $\omega$.
\begin{center}
\scriptsize
\begin{tabular}{@{}c r r r r r r@{}}
\toprule
$\omega$ & $k$ & $D$ & $A$ & $\log_2 W_{\rm est}$
& $\log_2 M_{\rm sym}^{\rm proxy}$ & status\\
\midrule
$2.81$ & 9 & 12 & $99{,}219{,}059$ & 128.833 & 93.337
& unit-leading estimate\\
$3$    & 9 & 12 & $99{,}219{,}059$ & 133.667 & 93.337
& unit-leading estimate\\
\bottomrule
\end{tabular}
\end{center}
Here the unit-leading dense proxy for symbolic preprocessing storage is
\begin{equation}
 M_{\rm sym}^{\rm proxy}=5B^2\binom{k+D}{D}\quad\text{bits}.
 \label{eq:pxl-symbolic-storage}
\end{equation}
This proxy assigns five bits to every stored base-field entry in the dense
symbolic array counted by the model.  It omits representation overhead,
temporary arrays, and storage traffic.  It is not the dense state of the
whole algorithm or a peak-memory bound.

The five charged Level-I terms have the following base-two logarithms.  The
terms depending on $\omega$ are recomputed in each row.
\begin{center}
\scriptsize
\begin{tabular}{@{}c r r r r r@{}}
\toprule
$\omega$ & $\log_2 C_1$ & $\log_2 C_2$ & $\log_2 C_3$
& $\log_2 C_{\rm fix}$ & $\log_2 C_{\rm lin}$\\
\midrule
$2.81$ & 102.354 & 84.399 & 109.027 & 109.525 & 112.877\\
$3$    & 109.275 & 84.399 & 109.027 & 109.525 & 117.924\\
\bottomrule
\end{tabular}
\end{center}
The work exponent uses the certified lower bound $p_{\rm dir}^{\rm L1}$ from
\eqref{eq:pxl-pdir-level1}, including the $19^{-4}$ transversality loss.
The augmented expression remains below the Level-I reference exponent $143$
whenever
\begin{equation}
 C_{\rm aux}
 <p_{\rm dir}^{\rm L1}2^{143}
  -150C_{\rm PXL}-C_{\rm target}.
 \label{eq:pxl-aux-threshold}
\end{equation}
Every returned candidate is checked against the complete residual system,
the signature format rule, and the original verifier.  The algorithm
therefore never returns an invalid signature.  Its completeness and displayed
cost estimate remain conditional on $\mathsf H_{\rm route}$ and hence on the
instancewise PXL rank predictions.  The affine support statement in
\Cref{prop:pxl-l2-affine-support} is exact and is not a solver premise.
\end{theorem}

\begin{proof}
Each trial samples a fresh retained target and an independent uniform affine
coset.
The exact target test is performed before solving.  By
\Cref{thm:sq-accepted-root}, one affine-coset trial contains an accepted root
with full column rank with probability at least
$p_{\rm reg}^{\rm L1}$.  Subtracting the exact nontransverse probability gives
$p_{\rm dir}^{\rm L1}$ as a certified lower bound.  Under
$\mathsf H_{\rm route}$, the actual matrices satisfy the published PXL model
and the added root extraction step recovers the required base-field root.  The
published algebraic work is the sum of the five costs in
\Cref{eq:pxl-C1,eq:pxl-C2,eq:pxl-C3,eq:pxl-Cfix,eq:pxl-Clin}.  The factor
$19^k$ already charges all assignments to the guessed variables, so no further
guessing multiplier is introduced.  Independent target and slice trials
contribute at most the reciprocal of $p_{\rm dir}^{\rm L1}$ in the modeled
expectation.  The target and setup
costs therefore belong inside the retry numerator.  The published-model
estimate charges the target term,
which is numerically smaller than the nonlinear solve, but leaves the other
costs, including the added root extraction work, in $C_{\rm aux}$.  Charging
each field operation at $150$ gates gives
the conditional expressions
\Cref{eq:pxl-estimated-gates,eq:pxl-total-gates}.  The final exact checks give
one-sided validity.
\end{proof}

The same direct $\ell=2$ model has more room at Levels III and V.  The
arithmetically checked estimates for $\omega=2.81$ are as follows.
\begin{center}
\scriptsize
\setlength{\tabcolsep}{4pt}
\begin{tabular}{@{}c r r r r r r r@{}}
\toprule
Level & $K$ & $h$ & $k$ & $D$ & $A$
& $\log_2 W_{\rm est}$ & $\log_2 M_{\rm sym}^{\rm proxy}$\\
\midrule
I   & 48 & 46 & 9  & 12 & $99{,}219{,}059$
    & 128.833 & 93.337\\
III & 72 & 70 & 13 & 17 & $2{,}710{,}186{,}080{,}528$
    & 187.030 & 137.856\\
V   & 96 & 94 & 15 & 23 & $547{,}365{,}920{,}911{,}569{,}552$
    & 245.042 & 186.447\\
\bottomrule
\end{tabular}
\end{center}
For the dense linear algebra model with $\omega=3$, the minimizing parameters
differ at Levels III and V.  They are therefore reported in a separate table
of estimates.
\begin{center}
\scriptsize
\setlength{\tabcolsep}{4pt}
\begin{tabular}{@{}c r r r r r r r@{}}
\toprule
Level & $K$ & $h$ & $k$ & $D$ & $A$
& $\log_2 W_{\rm est}$ & $\log_2 M_{\rm sym}^{\rm proxy}$\\
\midrule
I   & 48 & 46 & 9  & 12 & $99{,}219{,}059$
    & 133.667 & 93.337\\
III & 72 & 70 & 14 & 16 & $962{,}082{,}316{,}942$
    & 194.626 & 133.178\\
V   & 96 & 94 & 19 & 20 & $9{,}230{,}629{,}773{,}094{,}340$
    & 255.550 & 172.754\\
\bottomrule
\end{tabular}
\end{center}
Both storage columns report the unit-leading dense proxy
$M_{\rm sym}^{\rm proxy}$ from \eqref{eq:pxl-symbolic-storage}.  They do not
report total state or peak implementation memory.
The Level-III and Level-V rows use the same exact quotient, accepted-root
theorem, and conditional PXL model with their official dimensions.  They do
not supply evidence for the corresponding instancewise
$\mathsf H_{\rm route}$ premises beyond the Level-I analysis.

\subsection{Logical status}

The PXL estimate does not change the logical status of the public reduction.
The statements used by the Level-I route have the following status.

\begin{center}
\scriptsize
\begin{tabularx}{0.98\textwidth}{@{}>{\raggedright\arraybackslash}p{0.29\textwidth}>{\raggedright\arraybackslash}p{0.20\textwidth}X@{}}
\toprule
Part & Status & Meaning\\
\midrule
Zero-offset quotient and output split
& Exact after a public rank check
& Invertible public row reduction preserves every verifier equation.\\
Target consistency probability
& Exact
& A uniform target passes with probability $19^{-16}$, and the residual target is uniform after conditioning.\\
Affine-coset distribution
& Exact
& \Cref{lem:pxl-uniform-coset} samples every coset with the same probability.\\
Transversality incidence
& Exact
& An $A$-independent $h=v-2$ direction has nontransverse coset probability $19^{-4}$.\\
$\ell=2$ affine coefficient support
& Exact algebraic statement
& \Cref{prop:pxl-l2-affine-support} proves surjectivity, not uniformity.\\
Accepted full column rank root bound
& Exact under the stated spectrum inequality
& \Cref{thm:sq-accepted-root} and the transversality subtraction give the lower bound $p_{\rm dir}^{\rm L1}$ without an independence assumption.\\
Spectrum inequality
& Ideal transcript distribution or fixed-key certificate
& The byte-reduction transcript model gives a tail bound.  A fixed key can instead supply the projective rank certificate.\\
PXL solving degree and matrix ranks
& Ordinary MQ heuristic
& $\mathsf H_{\rm PXL}$ is the same kind of semi-regularity and rank model used in the PXL literature.\\
Application to each quotient instance
& Instancewise route premise
& $\mathsf H_{\rm route}$ states that the actual matrices follow the PXL model and that the added extraction step recovers an accepted root when one exists.\\
$\ell=4$ fixed-map support criterion
& Exact over the splitting field
& \Cref{thm:pxl-fixed-map-surjectivity} gives surjectivity exactly when every oil-quotient map is injective.  Descent requires base-field data.\\
$\ell=4$ coefficient family
& Structured-family heuristic, pending
& Ordinary MQ support does not cover the four Frobenius eigenblocks.  No $\ell=4$ row is released in this fragment.\\
PXL operation formula
& Unit-leading asymptotic estimate
& The expression charges all five PXL terms, target generation, and retries using the certified lower bound $p_{\rm dir}^{\rm L1}$.  Auxiliary costs remain explicit and symbolic.\\
Symbolic storage
& Unit-leading dense proxy
& $M_{\rm sym}^{\rm proxy}$ counts one dense symbolic array at five bits per base-field entry.  It is not a peak-memory bound.\\
Candidate validity
& Exact
& The complete residual system and original verifier check every output.\\
\bottomrule
\end{tabularx}
\end{center}

\subsection{Pending $\ell=4$ structured-family result}

% TODO(constrained-ledger): Replace this subsection only after the
% structured-family proof and the constrained ledger are complete.
% Do not infer an ell=4 result from the ordinary-MQ ell=2 calculation.
The following fixed-map statement isolates the exact injectivity condition
behind the split coefficient calculation.  It does not supply a numerical
$\ell=4$ row.

\begin{theorem}[Fixed-map channel surjectivity]
\label{thm:pxl-fixed-map-surjectivity}
Let $\Omega$ be a splitting field of odd characteristic.  Write
\[
 V=\bigoplus_{a=0}^{3}V_a,
 \qquad
 O=\bigoplus_{a=0}^{3}O_a,
\]
and fix linear maps $T_a:X\to V_a$.  Let
\[
 \overline T_a:X\longrightarrow V_a/O_a
\]
be the induced maps.  For each diagonal block, allow every symmetric form on
$V_a$ that vanishes on $O_a\times O_a$.  For each pair $a<b$, allow every
bilinear form on $V_a\times V_b$ that vanishes on
$O_a\times O_b$.  Pullback along the maps $T_a$ gives a linear map
$\Phi_T$ to the four diagonal and six cross quadratic forms on $X$.
Then
\[
 \Phi_T:\mathcal U_O\longrightarrow
 \operatorname{Sym}^2(X^*)^{10}
\]
is surjective if and only if every $\overline T_a$ is injective.  Here
$\mathcal U_O$ denotes the direct sum of the allowed diagonal and cross
coefficient blocks.
\end{theorem}

\begin{proof}
Suppose first that every $\overline T_a$ is injective.  Then $T_a(X)$ meets
$O_a$ only at zero.  On the direct sum $T_a(X)\oplus O_a$, prescribe any
symmetric form on $T_a(X)$ and prescribe zero on $O_a\times O_a$.  The
remaining mixed block can be chosen arbitrarily.  Extending this form to
$V_a$ realizes any diagonal quadratic form on $X$.

For a pair $a<b$, prescribe a bilinear form on
$T_a(X)\times T_b(X)$ whose value on $(T_a x,T_b x)$ is any chosen quadratic
form in $x$.  Such a bilinear lift exists because two is invertible.  Extend
it to
$(T_a(X)\oplus O_a)\times(T_b(X)\oplus O_b)$ while setting its block on
$O_a\times O_b$ to zero, and then extend it to $V_a\times V_b$.  The ten
blocks are independent, so these choices prove surjectivity.

Conversely, suppose that $\overline T_a$ has a nonzero kernel vector $x$.
Then $T_a x\in O_a$.  Every allowed diagonal form has zero value at
$(T_a x,T_a x)$.  The corresponding diagonal output therefore cannot equal
a quadratic form on $X$ that is nonzero at $x$.  Thus $\Phi_T$ is not
surjective.
\end{proof}

If all spaces, maps, and coefficient conditions in
\Cref{thm:pxl-fixed-map-surjectivity} arise by scalar extension from
$\F_{19}$, surjectivity over $\Omega$ descends.  Indeed, the cokernel becomes
zero after tensoring with $\Omega$, and faithful flatness shows that it was
already zero over $\F_{19}$.  This conclusion does not apply merely because
an eigenblock decomposition exists over $\Omega$.  The data and the map
$\Phi_T$ must themselves arise by scalar extension from the base field.

The $\ell=4$ shapes require a separate coefficient-family statement.  Their
affine reduction and spectrum theorem remain exact, but the source used to
justify the PXL parameters must respect the four Frobenius eigenblocks.  The
fixed-map theorem gives an exact support criterion after splitting.  It does
not show that the constrained official source meets that criterion after
variable fixing, or that the resulting systems have the selected degree and
Macaulay ranks.

No $\ell=4$ numerical row is part of this temporary replacement.  The final
text must wait for the constrained ledger for both official shapes.  Any
later row must use separate minimizing
parameters for each value of $\omega$ and must label its storage quantity as
a unit-leading dense proxy for symbolic preprocessing storage rather than
peak memory.
